\section{The expansion hierarchy}\label{sec:8.5}

The phonon-assisted absorption of silicon, \cref{eq:6.5-eps2PA}, accounts for the measured indirect edge so well that the phonon energies that are read off its thresholds agree with the neutron values to \SI{2}{\percent} (\cref{sec:6.5}), and yet it is a truncated formula. It keeps the electron--phonon vertex to second order and drops the fourth. It takes the phonons as harmonic, the electrons as independent particles, the thermal occupations to first order, and the nuclei as slow. Each of these choices drops the next term of an expansion in some small number. In a silicon detector at \SI{10}{mK} each small number has a value: the Born--Oppenheimer parameter is 0.067, the anharmonic parameter 0.049, the electron--phonon parameter at the edge $\sim10^{-2}$, the ratio of temperature to phonon energy $\sim10^{-5}$. Every rate in this book is built in the same way, and so every rate can be labeled by the orders it keeps. The label says where the next correction comes from and, with the small parameter, roughly how large it is.

\subsection{Small parameters}

\emph{Nuclear mass.} The Born--Oppenheimer parameter $\kappa_{\mathrm{BO}} = (m_e/M)^{1/4}$ of \cref{sec:5.1} is 0.0665 for natural silicon (\cref{we:5-kappa}). Through the hierarchy \cref{eq:5.1-scales} it sets the zero-point amplitude in units of the Bohr radius, the phonon energy in units of $\alpha^2m_e$ (as $\kappa_{\mathrm{BO}}^2$), the anharmonic level shifts ($\kappa_{\mathrm{BO}}^4$) and the nonadiabatic ones ($\kappa_{\mathrm{BO}}^6$). The Born--Oppenheimer approximation drops the non-adiabatic couplings $\mathcal{N}_{ji}$ of \cref{eq:5.1-N}, whose first-derivative and second-derivative terms are both built from the derivative coupling $\vect{d}^{\vect{R}\kappa}_{ji}$ of \cref{eq:6.1-d} (\cref{eq:6.1-Pd}). The effect of the non-adiabatic couplings on a given process depends on how fast the nuclei move. In zero-point motion a nucleus mixes excited electronic states into the ground state with an amplitude $\kappa_{\mathrm{BO}}^3 = 2.9\times10^{-4}$, and in an insulator the ratio that controls the mixing is $\omega_{\mathrm{ph}}/\Delta E$, which is $\omega_{\mathrm{opt}}/\Eg = 0.055 \approx 12\,\kappa_{\mathrm{BO}}^2$ in silicon (\cref{sec:5.1,sec:6.1}). A nucleus kicked by dark matter moves far faster, and the same coupling then produces the sudden amplitude \cref{eq:6.1-sudden1}, which is the Migdal effect.

\emph{Electron--phonon coupling.} A second-order amplitude is smaller than a first-order one by the coupling at the added vertex, divided by the off-shell distance of the intermediate state (\cref{sec:2.4}). For the electron--phonon coupling the square of that ratio is the \emph{off-shell parameter}
\begin{equation}
  \theta_{\mathrm{ep}} = \frac{\langle|g|^2\rangle}{\delta_{\mathrm{off}}^2} ,
  \label{eq:8.5-theta}
\end{equation}
where $\delta_{\mathrm{off}}$ is the distance of the intermediate electronic state from the energy shell and $\langle|g|^2\rangle$ is the mean squared vertex $g_{mn\nu}(\vect{k},\vect{Q})$ of \cref{eq:6.2-g} over the phonon modes and intermediate states that the amplitude visits. Each pair of vertices added to a rate costs one power of $\theta_{\mathrm{ep}}$. In silicon the optical vertices of the holes are about \SI{0.25}{eV} (\cref{sec:6.2}), and at the absorption edge the intermediate states lie 2 to \SI{2.5}{eV} off shell (\cref{sec:6.5}). These values give $\theta_{\mathrm{ep}} = 0.010$ to 0.016, the ``$\sim10^{-2}$ per mode'' that \cref{sec:6.5} estimated for the phonon-assisted amplitude. Acoustic vertices near the zone center are ten times smaller (\SI{26}{meV} at $Q = \SI{0.01}{\angstrom^{-1}}$) and their $\theta_{\mathrm{ep}}$ a hundred times smaller still. The parameter is not small everywhere. When an intermediate state approaches the shell, at the direct gap for light or near the edge $E_{\mathrm{min}}(q)$ for dark-matter scattering, $\delta_{\mathrm{off}}$ falls toward $|g|$ and $\theta_{\mathrm{ep}}$ reaches one at $\delta_{\mathrm{off}} \approx \SI{0.25}{eV}$. Closer still the expansion fails, the widths of \cref{eq:2.4-width} enter the energy denominators, and the resonant part becomes a cascade of real transitions (\cref{eq:2.4-split}). Away from the shell, where $\theta_{\mathrm{ep}} \ll 1$, lies the off-shell window in which \cref{sec:8.6} separates the zero-phonon from the phonon-dressed final states. Its edge is not where $\theta_{\mathrm{ep}}$ reaches one: at the margin of \SI{0.3}{eV} that \cref{sec:13.6} fixes for silicon, $\theta_{\mathrm{ep}} \approx 0.7$ bounds the correction from the next pair of vertices, whose size there is about 0.2.

\emph{Anharmonicity.} The anharmonic parameter $\eta_{\mathrm{anh}} = |\Phi^{(3)}|\ell_{\mathrm{zp}}/3\Phi^{(2)}$ of \cref{eq:5.7-eta} measures the zero-point amplitude against the length over which the force constants change, $\eta_{\mathrm{anh}} \sim \langle u\rangle/a$. By the counting of \cref{sec:5.7} it is of the size of $\kappa_{\mathrm{BO}}$, and in silicon it is 0.049. The cubic force constants are of first order in it and the quartic ones of second, so the three-phonon width and the first-order quartic shift are both of relative size $\eta_{\mathrm{anh}}^2 = 2.4\times10^{-3}$ (\cref{eq:5.7-order}). The silicon Raman line confirms this size: its measured width is $\gamma/\omega_{\mathrm{opt}} = 2.8\times10^{-3}$.

\emph{Temperature.} A detector at \SI{10}{mK} has $T = \SI{0.862}{\micro eV}$. The softest phonon that assists the silicon edge is the transverse acoustic phonon at X, $\omega_{\mathrm{TA}} = \SI{18.57}{meV}$. Against this phonon, $T/\omega_{\vect{Q}\nu} = 4.6\times10^{-5}$ and $n_{\vect{Q}\nu} = \ee^{-21550}$; against the gap, $T/\Eg = 7.4\times10^{-7}$. The ratio $T/\omega_{\vect{Q}\nu}$ is not small for every mode, since acoustic phonons near the zone center have energies below $T$. These phonons are, however, a fraction $(T/\omega_{\mathrm{D}})^3 = 4\times10^{-15}$ of the modes, with $\omega_{\mathrm{D}} = \SI{55.6}{meV}$ (\cref{we:5-debye}), and the expansion is in the occupations $n_{\vect{Q}\nu} = n_{\mathrm{B}}(\omega_{\vect{Q}\nu})$ that weight each mode, not in $T$ itself. The occupations enter every thermal formula of \cref{sec:2.5}, and they matter wherever a measured input was taken warm: the absorption edge of \cref{we:6-siedge} at \SI{300}{K}, the gap at room temperature, the phonon widths at \SI{77}{K}.

\emph{Recoil.} A nucleus given the momentum $\vect{q}$ would recoil freely with the energy $E_R = q^2/2M_\kappa$, and the \emph{recoil parameter}
\begin{equation}
  x_R = \frac{q^2}{2M_\kappa\langle\omega\rangle} = \frac{E_R}{\langle\omega\rangle} ,
  \label{eq:8.5-xR}
\end{equation}
compares that energy with the mean phonon energy $\langle\omega\rangle$ of the atom. The mean is taken over the phonon density of states per atom, normalized to three,
\begin{equation*}
  \langle\omega^k\rangle = \frac13\int_0^\infty\dd\omega\;\omega^k\,\frac{g_{\mathrm{ph}}(\omega)}{n_{\mathrm{at}}} ,
\end{equation*}
and for the density of states of \cref{sec:5.2} it is $\langle\omega\rangle = \SI{40.9}{meV}$. At $q = \SI{10}{keV}$, $E_R = \SI{1.91}{meV}$ and $x_R = 0.047$; $x_R$ reaches one at \SI{46}{keV}. Together with the Debye--Waller exponent $2W_\kappa$ of \cref{eq:5.6-W}, $x_R$ decides whether a momentum transfer creates one phonon, several, or a free recoil.

\emph{Suddenness.} A nucleus accelerated over a time $\tau_{\mathrm{acc}}$ excites an electronic state of energy $\Delta E$ with the sudden amplitude when $\Delta E\,\tau_{\mathrm{acc}} \ll 1$, and with an amplitude that falls faster than any power of $\Delta E\,\tau_{\mathrm{acc}}$ when that product is large (\cref{sec:6.1}). The \emph{suddenness ratio} $\Delta E\,\tau_{\mathrm{acc}}$ is the small parameter of the sudden limit. A silicon nucleus is kicked in about \SI{4e-21}{s}, while $1/\Eg = \SI{0.56}{fs}$, so that $\Eg\tau_{\mathrm{acc}} = 7\times10^{-6}$ across the gap. The lattice pulls the nucleus back on the slower time $1/\omega_{\mathrm{ph}}$, and \cref{sec:6.1} found that the pull-back changes the excitation probability by the factor $[1 - (\omega_{\mathrm{ph}}/\Delta E)^2]^{-2}$, 1.006 for silicon: the deceleration enters through the square of the adiabatic ratio $\omega_{\mathrm{ph}}/\Delta E$ of the first paragraph.

\emph{Multipoles.} The density vertex $\ee^{\ii\vect{q}_e\cdot\vect{r}}$ that excites an electron can be expanded in powers of the \emph{multipole parameter} $\xi_D = q_e r_{\mathrm{orb}}$, the momentum given to the electron times the size $r_{\mathrm{orb}}$ of the orbital it leaves or enters. The first nonvanishing term between orthogonal states is the dipole term of \cref{sec:2.2}. Three dark-matter processes give the electrons momenta that span six orders of magnitude. In electron scattering $q_e = q$, and at the typical $q_0 = \alpha m_e = \SI{1.89}{\angstrom^{-1}}$ an orbital of \SI{1}{\angstrom} gives $\xi_D \approx 1.9$: there is no multipole expansion, and the crystal form factors of \cref{chap:10} keep the plane wave whole. In the Migdal effect the electrons receive only the boost $m_e v_N$, \SI{100}{eV} for a silicon nucleus kicked by \SI{1}{GeV} dark matter at the largest momentum transfer (\cref{sec:6.1}), and $\xi_D \approx 0.05$. An absorbed boson of mass $m_\phi$ brings $q_e = m_\phi v \approx 10^{-3}m_\phi$, and $\xi_D = 5\times10^{-6}$ at $m_\phi = \SI{10}{eV}$.

\Cref{tab:8-params} collects the seven parameters. Each is an estimate with a coefficient of order one in front: the counting of \cref{sec:5.7} found coefficients between a quarter and two in silicon. Each is also small for silicon at the conditions of a dark-matter search, with three exceptions: $\theta_{\mathrm{ep}}$ near resonances, $x_R$ at large momentum transfer, and $\xi_D$ in electron scattering. Where a parameter is of order one the corresponding expansion is not used, and the label below records that the labeled quantity is computed to all orders in that parameter.

\begin{table}[tbp]
  \caption{Small parameters of the expansion hierarchy, with their values for silicon in a detector at \SI{10}{mK}. Each controls one component of the truncation label $\mathfrak{T}$ (\cref{eq:8.5-label}). ``Introduced'' is the section where the physics first appears; ``Number'' is where it becomes a correction to a dark-matter rate. The values are estimates, with coefficients of order one. The inputs are $M = \SI{26.16}{GeV}$ (natural silicon) for $\kappa_{\mathrm{BO}}$, the optical hole vertex of \SI{0.25}{eV} and off-shell distances of 2 to \SI{2.5}{eV} at the absorption edge for $\theta_{\mathrm{ep}}$, the bond coordinate of \cref{sec:5.1} for $\eta_{\mathrm{anh}}$, the transverse acoustic phonon at X (\SI{18.57}{meV}) for $T/\omega_{\vect{Q}\nu}$, the density of states of \cref{sec:5.2} ($\langle\omega\rangle = \SI{40.9}{meV}$) and the natural isotope mixture for $x_R$, a kick of \SI{4e-21}{s} across the \SI{1.17}{eV} gap for $\Delta E\,\tau_{\mathrm{acc}}$, and an orbital of \SI{1}{\angstrom} for $\xi_D$.}
  \label{tab:8-params}
  \centering\small
  \begin{tabular}{@{}>{\raggedright\arraybackslash}p{0.11\linewidth}>{\raggedright\arraybackslash}p{0.19\linewidth}>{\raggedright\arraybackslash}p{0.21\linewidth}>{\raggedright\arraybackslash}p{0.12\linewidth}>{\raggedright\arraybackslash}p{0.12\linewidth}>{\raggedright\arraybackslash}p{0.10\linewidth}@{}}
    \toprule
    Parameter & Definition & Silicon & Component & Introduced & Number \\
    \midrule
    $\kappa_{\mathrm{BO}}$ & $(m_e/M)^{1/4}$ & 0.0665; $\omega_{\mathrm{opt}}/\Eg = 0.055$, so $(\omega_{\mathrm{opt}}/\Eg)^2 = 3\times10^{-3}$ per step of $n_B$ in a lattice at rest & $n_B$ & \cref{sec:5.1} & \cref{chap:12,chap:13} \\
    $\theta_{\mathrm{ep}}$ & $\langle|g|^2\rangle/\delta_{\mathrm{off}}^2$, \cref{eq:8.5-theta} & 0.010--0.016 at the edge; 1 at $\delta_{\mathrm{off}} \approx \SI{0.25}{eV}$ & $n_g$ & \cref{sec:6.2} & \cref{chap:13} \\
    $\eta_{\mathrm{anh}}$ & $|\Phi^{(3)}|\ell_{\mathrm{zp}}/3\Phi^{(2)}$, \cref{eq:5.7-eta} & 0.049 ($\eta_{\mathrm{anh}}^2 = 2.4\times10^{-3}$) & $n_A$ & \cref{sec:5.7} & \cref{sec:11.5} \\
    $T/\omega_{\vect{Q}\nu}$ & temperature over phonon energy & $4.6\times10^{-5}$; $n_{\vect{Q}\nu} = \ee^{-21550}$ & $n_T$ & \cref{sec:2.5} & \cref{sec:13.7} \\
    $x_R$ & $q^2/2M_\kappa\langle\omega\rangle$, \cref{eq:8.5-xR} & 0.047 at \SI{10}{keV}; 1 at \SI{46}{keV} & $\mathcal{K}$ & \cref{sec:5.6} & \cref{sec:11.3,sec:11.4} \\
    $\Delta E\,\tau_{\mathrm{acc}}$ & excitation energy times kick duration & $7\times10^{-6}$ & $\mathfrak{M}$ & \cref{sec:6.1} & \cref{chap:12} \\
    $\xi_D$ & $q_e r_{\mathrm{orb}}$ & 1.9 (electron scattering at $q_0$); 0.05 (Migdal); $5\times10^{-6}$ (absorption at \SI{10}{eV}) & $n_D$ & \cref{sec:2.2} & \cref{sec:10.6,sec:12.2} \\
    \bottomrule
  \end{tabular}
\end{table}

\subsection{Lattice regimes}

The multiphonon expansion \cref{eq:5.6-multi} has a simple statistical reading. The number of phonons is Poisson distributed with mean $2W_\kappa$, and each phonon carries an energy drawn from the single-phonon distribution $P_1(\omega)$. For this compound Poisson distribution the cumulants of the energy deposited are $2W_\kappa$ times the moments of $P_1$, $c_k = 2W_\kappa\int\dd\omega\,\omega^kP_1(\omega)$. At $T = 0$ and on a cubic site, $P_1(\omega) \propto g_{\mathrm{ph}}(\omega)/\omega$ for $\omega > 0$, and $2W_\kappa = E_R\langle\omega^{-1}\rangle$ from \cref{eq:5.6-W}. The factors of $\langle\omega^{-1}\rangle$ cancel, and
\begin{equation}
  c_k = E_R\,\langle\omega^{k-1}\rangle , \qquad c_1 = E_R , \qquad c_2 = E_R\langle\omega\rangle .
  \label{eq:8.5-cumulants}
\end{equation}
The mean deposit is the free recoil energy, the sum rule \cref{eq:5.6-sum}. The variance is $E_R\langle\omega\rangle$. The width of the spectrum relative to its mean is
\begin{equation}
  \frac{c_2^{1/2}}{c_1} = \Bigl(\frac{\langle\omega\rangle}{E_R}\Bigr)^{1/2} = x_R^{-1/2} .
  \label{eq:8.5-width}
\end{equation}
\Cref{eq:8.5-width} is the reason for defining $x_R$ with the arithmetic mean $\langle\omega\rangle$, while the Debye--Waller exponent carries the inverse mean: $2W_\kappa = x_R\,\langle\omega\rangle\langle\omega^{-1}\rangle \ge x_R$, with $\langle\omega\rangle\langle\omega^{-1}\rangle = 1.35$ in silicon ($\langle\omega^{-1}\rangle^{-1} = \SI{30.4}{meV}$, \cref{we:5-dw}). The exponent $2W_\kappa$ counts phonons; the recoil parameter says whether the deposit is concentrated near $E_R$.

Together, the exponent and the recoil parameter fix the lattice regime. For $2W_\kappa \ll 1$, and therefore $x_R \ll 1$, most scatterings leave the lattice unchanged and a fraction $2W_\kappa$ creates one phonon, which carries an energy of the order of $\langle\omega\rangle$, far above $E_R$. In this \emph{one-phonon} regime the coherent response of the crystal (\cref{sec:11.3}) decides the rate. In silicon the one-phonon regime holds below about \SI{10}{keV}, where $2W_\kappa \le 0.063$ and $x_R \le 0.047$. For $2W_\kappa$ of order one, several terms of \cref{eq:5.6-multi} compete and the deposit comes in a few whole phonons: the \emph{multiphonon} regime, with $2W_\kappa = 1$ at \SI{40}{keV} and $x_R = 1$ at \SI{46}{keV}. For $x_R \gg 1$ the spectrum narrows around $E_R$, and by the central limit theorem it approaches a Gaussian of width $(E_R\langle\omega\rangle)^{1/2}$: the free recoil, broadened by the motion of the nucleus in its well. In this \emph{impulse} regime the spectrum is the elastic recoil of \cref{sec:1.3}, with the width as its leading correction~\cite[Sect.~IV\,A, Eqs.~(90)--(92)]{kahn2022searches}. The width falls below $E_R/3$ for $x_R \ge 9$, at $q \ge \SI{139}{keV}$ or $E_R \ge \SI{0.37}{eV}$ in silicon. The Einstein oscillator of \cref{sec:1.3}, at \SI{60}{meV}, puts the same boundary at \SI{0.54}{eV}; the computed density of states has a lower mean phonon energy and moves the boundary down by a third.

Temperature acts differently on the exponent $2W_\kappa$ and on the recoil parameter. Thermal phonons add the factor $\coth(\omega/2T)$ to $2W_\kappa$ (\cref{eq:5.6-W}), raising the number of phonons absorbed and emitted, while the mean deposit stays at $E_R$ at any temperature. At \SI{10}{mK} the factor is one for every mode that matters, and the regimes are the same as at $T = 0$. The boundaries between the regimes are gradual, and a formula is placed in a regime by what it assumes, not by the momentum at which it is evaluated.

\subsection{The truncation label}

Each small parameter has an order to which a formula keeps it. Together with three labels that are not orders, the orders form the \emph{truncation label}
\begin{equation}
  \mathfrak{T} = (n_B,\, n_g,\, n_A,\, \mathcal{X},\, n_T,\, \mathcal{K},\, \mathfrak{M},\, n_D) .
  \label{eq:8.5-label}
\end{equation}
The five integers are the Born--Oppenheimer, electron--phonon, anharmonic, thermal and multipole orders. One rule fixes them all: \emph{an order is the power of its expansion parameter carried by the quantity being labeled}, counted in the response or the rate, not in the number of vertices. The rule generalizes the counting of \cref{sec:5.7}, where a width with two cubic vertices and a shift with one quartic vertex are of one order. The three remaining components are labels: the many-body approximation $\mathcal{X}$, the lattice regime $\mathcal{K}$ and the Migdal regime $\mathfrak{M}$. The Born--Oppenheimer order $n_B$ always appears with its value, as in $n_B = 0$; the value distinguishes it from the Bose function $n_{\mathrm{B}}(\omega)$.

\emph{Born--Oppenheimer order $n_B$}, the power of the derivative coupling $\vect{d}^{\vect{R}\kappa}_{ji}$ in the quantity. At $n_B = 0$ the electrons are in the adiabatic state of the instantaneous nuclear configuration. The order $n_B = 0$ covers the clamped-nucleus band structures of \cref{chap:3}, the response functions of \cref{chap:4}, and the phonons on the Born--Oppenheimer surface of \cref{chap:5}, all of which are evaluated without any derivative coupling. The sudden amplitude \cref{eq:6.1-sudden1} is linear in $\vect{d}$, so the probability $1 - P_0$ that a kick excites the electrons is at $n_B = 2$, and so is the Migdal rate. The odd order $n_B = 1$ appears only in interference, between an amplitude that contains the coupling and one that does not; \cref{sec:8.6} meets such a term. The cost of a step depends on how fast the nuclei move, because each power of $\vect{d}$ enters with the nuclear velocity, as $\vect{v}_N\cdot\vect{d}/\Delta E$ (\cref{eq:6.1-sudden1}). For a lattice in its vibrational ground state the nuclear velocity is the velocity of zero-point motion. The step $n_B \to n_B + 2$ then multiplies an adiabatic quantity by $1 + O\bigl((\omega_{\mathrm{ph}}/\Delta E)^2\bigr)$, and the correction is $3\times10^{-3}$ for the optical phonon across the gap of silicon. The amplitude $\kappa_{\mathrm{BO}}^3$ above is the same ratio in the atomic counting, with $\Delta E \sim \alpha^2m_e$. For a nucleus kicked by dark matter the nuclear velocity is $\vect{q}/M_\kappa$, and the step from $n_B = 0$ to $n_B = 2$ is the sudden probability itself. The electrons that follow a nucleus carry a share $\kappa_{\mathrm{BO}}^4 = m_e/M$ of an f-sum (\cref{sec:8.6}); this share is fixed by the masses alone and is not a step of this kind.

\emph{Electron--phonon order $n_g$}, the power of the vertex $g$. Direct absorption, \cref{eq:6.5-eps2dir}, is at $n_g = 0$; phonon-assisted absorption, \cref{eq:6.5-eps2PA}, has one vertex in the amplitude and is at $n_g = 2$, as are the one-phonon emission rates and the Fan--Migdal widths of \cref{sec:6.3}. The vertex itself is of first order in the displacement and is at $n_g = 1$. Raising $n_g$ by two costs one power of $\theta_{\mathrm{ep}}$.

\emph{Anharmonic order $n_A$}, the power of $\eta_{\mathrm{anh}}$. The harmonic crystal of \cref{sec:5.2,sec:5.3,sec:5.6} is at $n_A = 0$, and so is the multiphonon formula. The three-phonon width \cref{eq:5.7-gamma} and the first-order quartic shift are both at $n_A = 2$, and so is the leading anharmonic phonon self-energy; at $n_A = 1$ the energies do not change, since the average of the cubic term over a harmonic state vanishes. First-principles electron--phonon theory starts from the point $n_B = 0$, $n_A = 0$: its lattice Hamiltonian rests only on the harmonic and the Born--Oppenheimer approximations~\cite[Sect.~III\,A, text after Eq.~(19)]{giustino2017electron}.

\emph{Many-body approximation $\mathcal{X}$}, one of five: the independent-particle bubble \cref{eq:4.6-R0}, IP (the Lindhard function of \cref{sec:4.5} in the electron gas); the random-phase approximation (RPA) \cref{eq:4.6-rpa} with the head alone, or with local fields (LF) \cref{eq:4.6-rpaLF}, written RPA + LF; an exchange--correlation kernel of time-dependent density-functional theory (TDDFT) \cref{eq:4.6-tddft}; and the Bethe--Salpeter equation (BSE) \cref{eq:4.9-bse}. $\mathcal{X}$ is not an order. Density-functional perturbation theory supplies the phonons, Born charges and $\varepsilon_\infty$ as static self-consistent responses with the Hartree and exchange--correlation kernel of the functional, and these inputs carry $\mathcal{X}$ = TDDFT at $\omega = 0$, whatever the approximation of the electronic response they are combined with.

\emph{Thermal order $n_T$}, the power of the phonon occupations $n_{\vect{Q}\nu}$. At $n_T = 0$ there is no thermal occupation, though zero-point motion is kept wherever phonons enter; at $n_T = 1$ a rate is linear in the occupations, as \cref{eq:6.5-eps2PA} is through $n_{\vect{Q}\nu} + \frac12 \pm \frac12$. The thermal occupation of electrons across a gap, $\ee^{-\Eg/T}$, is smaller than any power and is dropped at every order.

\emph{Lattice regime $\mathcal{K}$}, one of one-phonon, multiphonon and impulse: the phonon content of the final state. For a coupling to the nuclei it is read from $x_R$ and $2W_\kappa$ as above. The single term $n = 1$ of \cref{eq:5.6-multi} is one-phonon; the full sum is multiphonon; its Gaussian limit around $E_R$ is impulse. A process is one-phonon as well when it creates its phonon through one electron--phonon vertex in the amplitude, as phonon-assisted absorption does. Without phonons in the final state the label does not apply.

\emph{Migdal regime $\mathfrak{M}$}, sudden or adiabatic. The boost \cref{eq:6.1-boost} and the hydrogen survival probability \cref{eq:6.1-P0} are sudden; the other value is the adiabatic limit of \cref{sec:6.1}, in which the electrons follow the nucleus and end in their ground state. The label applies to a formula that treats the struck nucleus as a classical trajectory. Where the phonon final states are kept, at small $x_R$, the same physics is carried by $n_B$ and $\mathcal{K}$ instead (\cref{sec:12.1}).

\emph{Multipole order $n_D$}, the power of $\xi_D$. At small momentum transfer the dipole transition $1s \to 2p$ of hydrogen has $|f|^2 \simeq 0.555\,q^2a_0^2$ (\cref{sec:2.2}) and is at $n_D = 2$, the lowest order at which a transition between orthogonal states occurs. The full plane wave, kept whole, has no multipole order.

The \emph{base point} of the label is
\begin{equation}
  \mathfrak{T}_0 = (n_B = 0,\, n_g = 0,\, n_A = 0,\, \mathcal{X} = \mathrm{IP},\, n_T = 0;\ \mathcal{K},\, \mathfrak{M},\, n_D\ \text{not applicable}) :
  \label{eq:8.5-base}
\end{equation}
clamped nuclei, no electron--phonon vertex, independent particles, zero temperature, and no expansion of the plane wave, with no phonon in the final state and no nuclear recoil to classify. The unscreened crystal form factor of \cref{chap:10} is at $\mathfrak{T}_0$. When a formula takes its phonons or Born charges from density-functional perturbation theory, it keeps $\mathcal{X}$ for its electronic response and names those inputs as static TDDFT. A formula elsewhere on the label is written by listing only the components that differ from $\mathfrak{T}_0$. The direct dielectric function \cref{eq:6.5-eps2dir} is $\mathfrak{T}_0$ in this notation, a three-phonon linewidth is ``$n_A = 2$'', and the screened loss function with local fields is ``$\mathcal{X} = \mathrm{RPA + LF}$''. Qualifiers that are not orders, such as ``unscreened'' or ``with spin--orbit coupling'', are added in words. \Cref{sec:9.4} places the literature formulas on the label.

Several approximations are deliberately left off the label, because none of them is an expansion in a small number with a next order that could be computed. Of this kind are the exchange--correlation functional, the frozen core of a pseudopotential, the $GW$ quasiparticle energies, and the numerical parameters of a calculation (grids, cutoffs, Wannier interpolation); their errors are measured by tests and convergence. Methods that sum an expansion to all orders lie off the finite values of an axis and are described in words. The self-consistent harmonic approximation of \cref{sec:5.7} resums the anharmonic expansion, molecular dynamics on the full Born--Oppenheimer surface is exact in $n_A$ at $n_B = 0$, and a classical cascade calculation is off the perturbative axes altogether.

\subsection{Many-body approximations as resummations}

The approximations of \cref{sec:4.6,sec:4.9} are ordered by what they include, and each is a limit of the next, but they are not successive orders in a small number. Each is a resummation. The random-phase approximation already sums the Coulomb interaction to all orders through the Dyson equation, the kernel of time-dependent density-functional theory adds a correction inside that sum, and the Bethe--Salpeter equation adds a ladder of electron--hole attractions. An expansion would need the kernel to be small against the Coulomb interaction it corrects, and the kernel is not small. In silicon the local fields lower $\varepsilon_\infty$ by \SI{11}{\percent}, from $11.45$ to $10.22$; the adiabatic kernel restores about half of this reduction, to $10.81$; and adding the attraction to the local fields raises the local-field value by \SI{17}{\percent}, to $\approx 11.9$ (\cref{we:4-siepsinf,sec:4.9}). Each correction moves $\varepsilon_\infty$ by about as much as the last, and not always in the same direction, so there is nothing to expand in. At the lower absorption peak the approximations disagree by a factor of more than two, while above \SI{4}{eV} they agree within the accuracy of the bands.

The error of a many-body approximation is therefore estimated by computing the next one, not by a power. The difference between two computed approximations is the many-body uncertainty of a rate: for the screening factor $1/\varepsilon_\infty^2$ of silicon it is up to a third across the hierarchy, and for light dark matter near the gap it is the leading uncertainty. \Cref{sec:8.8} turns these computed differences, and the next-order terms of the other components, into the theory uncertainty of a rate.

\begin{workedexample}{The label of phonon-assisted absorption}\label{we:8-label}
We read the label of the phonon-assisted dielectric function \cref{eq:6.5-eps2PA} off its assumptions, component by component, and estimate the size of each first omitted term in silicon at the indirect edge.

The phonon frequencies and vertices come from density-functional perturbation theory on the Born--Oppenheimer surface, with the electrons responding statically, so the formula is at $n_B = 0$, and these inputs are at $\mathcal{X}$ = TDDFT in its static limit. The static screening of the vertex misses corrections of relative order $(\omega_{\mathrm{ph}}/\Eg)^2$ (\cref{sec:6.2}), $3\times10^{-3}$ for the optical phonon. The phonons are harmonic, $n_A = 0$; their anharmonic widths, $\gamma/\omega = 2.8\times10^{-3}$ for the optical phonon, would broaden each threshold by a fraction of a meV.

The amplitude $M^\pm_{cv\nu}$ contains one optical vertex and one electron--phonon vertex, so the rate is quadratic in $g$: $n_g = 2$. The first omitted order, $n_g = 4$, adds processes with two phonons and the dressing of the one-phonon process by a second. It is smaller by $\theta_{\mathrm{ep}} \approx 10^{-2}$ below $E_{\mathrm{dir}}$, where every intermediate state lies electronvolts off shell, and it is not small near $E_{\mathrm{dir}}$, where the intermediate electron becomes real and \cref{sec:6.5} replaced the expansion by the width $\varGamma_{\mathrm{w}}$.

The occupation factor $n_{\vect{Q}\nu} + \frac12 \pm \frac12$ is linear in $n_{\vect{Q}\nu}$, so the formula is at $n_T = 1$. Linearity is exact at $n_g = 2$ for harmonic phonons: a term quadratic in the occupations needs two phonons and therefore $n_g = 4$. At \SI{10}{mK} the occupations are $\ee^{-21550}$ and smaller, and only the emission branch survives; the label still records $n_T = 1$, because the formula holds at the \SI{300}{K} and \SI{77}{K} of the measured edge in \cref{we:6-siedge}. The electrons are thermally unoccupied across the gap, $T \ll \Eg$, as at every order.

The bands are independent-particle bands with a rigid shift, and the electron and hole do not attract each other: $\mathcal{X} = \mathrm{IP}$. Each event creates or absorbs exactly one phonon, through the vertex $g$, so $\mathcal{K}$ = one-phonon; the photon vertex acts on the electrons and no nucleus recoils, so $\mathfrak{M}$ does not apply. The dielectric function $\varepsilon_2$ is defined in the limit $\vect{q} \to 0$, where no multipole expansion is made, so $n_D$ does not apply either. In full,
\begin{multline*}
  \mathfrak{T}\bigl[\varepsilon_2^{\mathrm{PA}}\bigr] = (n_B = 0,\, n_g = 2,\, n_A = 0,\, \mathcal{X} = \mathrm{IP},\, n_T = 1,\\
  \mathcal{K} = \text{one-phonon};\ \mathfrak{M},\, n_D\ \text{not applicable}) ,
\end{multline*}
written as ``$n_g = 2$, $n_T = 1$, $\mathcal{K}$ = one-phonon'' against the base point, with $\mathcal{X}$ at its base value IP. A calculation for a detector at millikelvin temperatures drops the absorption branch and so puts the same formula at $n_T = 0$; the phonon-assisted loss function of \cref{app:E.14} is labeled that way.

The label also ranks the corrections. The next electron--phonon order costs $10^{-2}$, the static screening of the vertex a few $10^{-3}$, and the anharmonic widths less. The excitonic correction, however, is not a percent effect at the edge. The electron--hole attraction moves every threshold down by the exciton binding energy, \SI{14.7}{meV} (\cref{sec:4.9}), comparable to the energy of the transverse acoustic phonon that assists the transition, and the measured kinks of \cref{we:6-siedge} are exciton thresholds. Near the edge, then, the largest error of \cref{eq:6.5-eps2PA} lies on the axis that has no small parameter, $\mathcal{X} = \mathrm{IP} \to \mathrm{BSE}$, and it is estimated by solving the Bethe--Salpeter equation.
\end{workedexample}
